\section{Commercial orchestration and the recovery of a causal batch}
\label{exec22:workflow-scope}

The chapter compares commercial orchestration, its open clients and delegates, and a distributed execution alternative. Each successive treatment makes the boundary more concrete: a manual describes an operating choice; a client exposes its request and return path; a delegate reveals how retries, state and dependencies are handled. The final Ray comparison asks when distributing the same scientific workload is worth the added recovery and accounting obligations. The evidence supports conditional reuse, with complete scientific outputs and all attempted requests as the common comparison.

A recovered causal batch is useful only if its result answers the declared
question for the declared population and information set. Consider a nightly
offer analysis that fails after writing some cohort effects. Restarting the job
can save computation, but it can also duplicate customers, reuse an outdated
estimate or execute a changed procedure. The relevant product requirement joins
execution recovery to preservation of the scientific request. This motivates
comparing available orchestration before building another scheduler.

The six selected official manuals cover Databricks Lakeflow Jobs on AWS,
Snowflake task graphs, Azure Machine Learning API-v2 pipeline reuse, Dataiku
DSS scenarios and AWS Step Functions. They provide current technical
observations retrieved on 22 October 2026. They do not establish an unchanged
capability at the literature cutoff, 7 October 2026. Databricks' repair and
monitoring pages display 9 and 15 October updates. Azure displays a 30 March
date while its HTML repository metadata records a 19 October update. The
remaining mutable manuals lack an authenticated historical revision. These
observations supplement the cutoff survey; historical coverage remains open.
\citep{exec22_workflow_databricks_repair,exec22_workflow_databricks_monitor,
exec22_workflow_azure_reuse,exec22_workflow_snowflake_graphs,
exec22_workflow_dataiku_scenarios,exec22_workflow_aws_errors}

\subsection{Partial repair must preserve both writes and meaning}

The inspected Databricks repair manual describes rerunning unsuccessful tasks
and their dependents using current job and task settings. Each unsuccessful
task starts again from the beginning. The text explicitly warns that Jobs does
not make task writes idempotent: a partially completed append can duplicate data
when repaired. A shared job cluster is recreated using current settings, and
repair parameters can override earlier values. These mechanisms can support
efficient recovery; they make the application's output-commit and version
contracts material to a causal result. They do not establish that the platform
cannot implement those contracts.~\citep{exec22_workflow_databricks_repair}

Let the finite target population be $\mathcal U$, and let $\tau_i$ be its
stipulated potential-outcome contrast $Y_i(1)-Y_i(0)$. The equal-unit target is
\begin{equation}
 \tau_{\mathcal U}=\frac{1}{|\mathcal U|}\sum_{i\in\mathcal U}\tau_i.
 \label{exec22:workflow-target}
\end{equation}
For an exact constructed example, three units have contrast 6 and one has
contrast $-2$, in hypothetical dollars per unit. Their effect is 4. If the
first two positive-contrast units are written before interruption, and all
four units are appended on repair, the stored-row average becomes
\begin{equation}
 \tau_{\rm append}=\frac{5(6)+1(-2)}{6}=\frac{14}{3},
 \qquad \tau_{\rm append}-\tau_{\mathcal U}=\frac{2}{3}.
 \label{exec22:workflow-duplicates}
\end{equation}
The batch is computationally complete but its implicit positive-group weight
has changed from $3/4$ to $5/6$. At stipulated per-unit cost 1, the intended
net value 3 becomes the misleading value $11/3$. No statistical noise is
needed for this error. These are exact authored values, rather than bank
observations or measured vendor behavior.~\citep{exec22_workflow_reference}

A final output keyed by unit identifier can recover four unique units, provided
the values belong to the same semantic request and the final commit is atomic
or otherwise conflict checked. The history still contains six write attempts.
Deduplicating the result must not erase the partial attempt, its cost or its
failure. Conversely, overwriting a key with a result from a different target
or data vintage is not valid idempotency. Identifier equality alone does not
establish scientific equality. The proposed adapter must distinguish a repeated
delivery of the same result from a conflicting new result.

\subsection{Job status and request status answer different questions}

The current Databricks monitoring manual distinguishes ``Succeeded with
failures'' from ordinary success and explains job outcomes in terms of leaf
tasks. A final task that tolerates an upstream failure can complete without
establishing that every requested customer cohort or sensitivity scenario was
evaluated. A causal batch therefore needs a request-level ledger alongside the
scheduler's result: completed, refused, failed, skipped and reused requests
remain distinguishable.~\citep{exec22_workflow_databricks_monitor}

The same manual describes a 48-hour finished-run graph and a 60-day run-history
window. It also explains that one-time API submissions have no saved job name,
so searching by that name does not find them. A durable scientific record
should bind the actual run identifier and export the relevant lineage before
its operating display expires. Retention and search convenience are useful
platform features; neither determines the validity of the causal assumptions
or the completeness of the requested workload.
\citep{exec22_workflow_databricks_monitor}

For a declared request set $\mathcal Q$, report the requested count together
with each disposition count. Accuracy among returned results and the fraction
of requests receiving an admissible result are different quantities. Retry
success must not silently remove difficult requests from the denominator.
A request outside the method's supported intervention domain needs an
unsupported scientific disposition; a transient storage failure is an operating
failure. A refusal because declared premises conflict is distinct from both. Treating them identically would hide the
appropriate remedy and distort the comparison with another route.

\subsection{Version the scientific dependencies behind an execution DAG}

Snowflake's inspected task-graph manual describes dependent SQL execution,
retries, finalizers and overlap controls. Its graph is a scheduling DAG. An
edge that says one job must finish before another does not imply a causal
arrow or identify an intervention. This distinction matters when an agent
translates the same word ``graph'' across workflow and causal APIs.
\citep{exec22_workflow_snowflake_graphs}

The manual says that graph versions include task properties, but not the
definitions of stored procedures called by those tasks. If a called procedure
changes during execution, its new definition can run within the current graph
run. Scientific replay must therefore pin procedure definitions and their
relevant dependencies separately. A stable task-graph version is insufficient
to infer unchanged estimator code.
\citep{exec22_workflow_snowflake_graphs}

The same text says that a suspended child can be treated as though it
succeeded, while a finalizer does not start when the root task is skipped.
Overlap can be allowed where read/write operations tolerate it. These are
documented execution choices whose compatibility with the proposed experiment
must be assessed. A manifest and commit protocol should make incompatible
concurrent writes detectable and should retain skipped scientific requests.
This is a design requirement, not a demonstrated Snowflake adapter or an
absence claim about other product features.
\citep{exec22_workflow_snowflake_graphs}

\subsection{A reusable output must still answer the same target}

Azure Machine Learning's inspected API-v2 description requires a deterministic
component, ForceReRun=false and matching code, environment, inputs, parameters,
output settings and run settings for reuse. The reused component displays the
earlier status, output, logs and metrics. Force-rerun children cannot be reused
by other jobs. These controls offer a useful starting point for lineage and
reproducibility, with the historical-edition limitation stated above.
\citep{exec22_workflow_azure_reuse}

The scientific application must express its dependencies through the declared
artifacts. An unchanged URI need not establish an unchanged table behind it;
an undeclared population revision can change the question even if the code
is identical. The determinism flag is a declaration about the component, not
an independent validation of determinism or causal assumptions. The manual's
code-snapshot and environment comparisons help diagnose execution reuse; the
library must add the estimand and information clock to the meaning of reuse.
\citep{exec22_workflow_azure_reuse}

In the finite example, keeping group effects $(6,-2)$ but changing population
weights from $(3/4,1/4)$ to $(1/4,3/4)$ changes the effect from 4 to 0 and net
value from 3 to $-1$. Retaining the original weights but revising the last
unit's contrast from $-2$ to $-10$ changes the effect to 2 and net value to 1.
The first change is a target change; the second is a data-vintage change.
A cached number may be exactly equal to an earlier output and still answer
the wrong current question.~\citep{exec22_workflow_reference}

\subsection{Retries need a budget that survives redrive}

The current AWS Step Functions manual describes ordered retriers and catchers,
backoff, jitter and bounded retry counts. Catchers apply to Task, Parallel and
Map states, rather than top-level execution failures. Eligible states can have
their retry counts reset on redrive. A library request that requires a total
budget must retain attempts across these boundaries, as well as configure each
state's retry policy. A new execution identifier must not erase earlier time,
resource use or failed scientific requests.
\citep{exec22_workflow_aws_errors}

Error classes require the exact documented meaning. In the inspected current
text, the States.ALL wildcard excludes DataLimitExceeded and Runtime. The
manual explicitly permits a named DataLimitExceeded retry or catch, whereas
Runtime remains non-retryable. It would be inaccurate to infer from the
wildcard exclusion that every explicit payload-limit handler is unavailable.
The historical revision must be recovered before attributing these semantics
to the cutoff edition.~\citep{exec22_workflow_aws_errors}

The manual states that retries are billed as state transitions. Actual
comparison requires more than the number of successful final tasks. For a
common workload, define a proposed accounting identity
\begin{equation}
 C_{\rm workflow}=C_{\rm transitions}+C_{\rm compute}+C_{\rm storage}
                  +C_{\rm transfer}+C_{\rm human}+C_{\rm license}.
 \label{exec22:workflow-cost}
\end{equation}
Each term includes failed and recovery attempts within the declared horizon;
overlapping charges must not be counted twice. Rates, region, service edition,
retention and staff time need authentic evidence before evaluating this
identity numerically. No observed prices, runtime advantage or bank return
are supplied by the exact arithmetic example.

\subsection{What remains unresolved in the commercial comparison}

Dataiku's retrieved page identifies DSS 15 and indexes automation scenarios,
build/check/metrics/code steps, trigger categories and failure/conditional
execution. It is primarily an index. Those headings establish the scope of
linked topics, not detailed transaction, retry or concurrency semantics.
Section~\ref{exec23:dss-scope} develops the technical step, trigger and test controls from separately saved manuals. That resolves the index-only limitation for those mechanisms, while the exact historical server-patch identity remains unresolved.~\citep{exec22_workflow_dataiku_scenarios}

SAS Viya orchestration, Domino and Posit Workbench/Connect receive selected open-client and delegate comparisons in Sections~\ref{exec23:clients-scope} and~\ref{sec:exec24-workflow-contracts}. Those source findings resolve part of this initial candidate inventory; they do not qualify licensed servers, complete capability coverage or supported deployment combinations. Complementary market discovery, product aliases, deployment editions,
licensing and actual prices remain necessary. For the five inspected families,
current manuals support specific mechanisms and application obligations;
none completes historical coverage or native qualification. This comparison
therefore motivates an architecture that reuses suitable orchestration while
adding scientific semantics. Whether each adapter preserves those semantics
remains a proposition to test.

\subsection{A common manifest across the library, CLI and MCP}

Let $m(q)$ serialize the full declared meaning of scientific request $q$:
\begin{equation}
 m(q)=(D,T,I,G,A,E,S,N,O),\qquad k(q)=\operatorname{hash}(m(q)).
 \label{exec22:workflow-manifest}
\end{equation}
Here $D$ identifies input snapshots and availability, $T$ the population and
estimand, $I$ the information clock, $G$ graph/identification assumptions,
$A$ the estimator and covariance settings, $E$ code and environment, $S$ seeds
or randomness coupling, $N$ numerical tolerances, and $O$ the output schema.
The serialization must be canonical and must include the dependencies needed
for that method; the tuple is a proposed contract, not a proof of completeness.
Equal hashes bind specified bytes and declarations. They do not establish
truthful provenance, valid assumptions or an identified causal effect.

The reference program constructs a finite manifest and verifies that changes
to the data, target weights or environment change its key. It also verifies
the append bias and keyed recovery using exact rational arithmetic. This is
a local authored reference, not a vendor adapter or a fault-injected deployment.
The CLI and MCP should transport the same manifest and return lineage, the
target, uncertainty meaning and request disposition. A refused unidentified
effect must not be converted into an ordinary numerical answer by either
entry point.~\citep{exec22_workflow_reference}

\subsection{Demonstrate merit on public fixtures before a costly trial}

A first frozen experiment should combine a public dataset with a separately
generated potential-outcome fixture. The public source exercises realistic
schema, scale and custody; generated outcomes supply known causal truth under
declared assumptions. Public observational agreement alone supplies no such
truth label. Compare uninterrupted and interrupted execution with the same
targets, data, methods, randomization and resource envelope.

Prespecify partial writes, lost responses after commit, changed tables or
procedures, stale cache, skipped children, payload limits and exhausted redrive
budgets. Every requested case remains in the denominator. Deterministic exact
cases require identical results and statuses. Floating or randomized cases
need predeclared tolerances and appropriate coupled randomness. Incompatible
manifests must be refused; commits must be unique and the attempt/cost ledger
complete. A statistical estimator retains a separate calibration protocol.
Computational equivalence does not itself establish coverage or identification.

This experiment can justify whether a native trial is worth its cost. Later
commercial qualification needs the chosen edition, licensed environment,
actual execution and complete prices; later bank acceptance needs authentic
bank questions, truth and owners. The proposed contract supplies a concrete
route for those stages while preserving uncertainty about their outcomes.
It does not establish that any sibling repository or commercial platform has
already adopted the causal library.
